\subsection{ব্যাংকের সম্পর্কিত টেবিলে SQL-এর প্রয়োগ}

ব্যাংকের ডেটাবেজ একটি গুরুত্বপূর্ণ বাস্তব উদাহরণ। এখানে একই সঙ্গে
শর্ত, সমষ্টিগত ফাংশন, UPDATE এবং সম্পর্কিত টেবিলের কুয়েরি
বোঝা যায়।

\begin{chSixSideBySide}[.48]
\subsubsection{BankCustomer Table}



\small

\begin{tabular}{|l|l|l|l|}
\hline

\chSixHeader
CustomerID & Name & Address & Phone\\
\hline
C01 & Rafi & Rajshahi & 01712000001\\\hline
C02 & Nila & Dhaka & 01712000002\\\hline
C03 & Sami & Rajshahi & 01712000003\\\hline
C04 & Mitu & Khulna & 01712000004\\
\hline

\end{tabular}
\chSixNextColumn
\subsubsection{Account Table}



\small
\setlength{\tabcolsep}{2.5pt}

\begin{tabular}{|l|l|l|l|}
\hline

\chSixHeader
AccountNo & CustomerID & AccountType & Balance\\
\hline
A1001 & C01 & Savings & 150000\\\hline
A1002 & C02 & Current & 85000\\\hline
A1003 & C01 & Current & 70000\\\hline
A1004 & C03 & Savings & 220000\\\hline
A1005 & C04 & Savings & 45000\\
\hline

\end{tabular}
\end{chSixSideBySide}
\begin{chSixSideBySide}[.48]
\subsubsection{Loan Table}



\small

\begin{tabular}{|l|l|l|l|}
\hline

\chSixHeader
LoanID & CustomerID & LoanType & Amount\\
\hline
L01 & C01 & Home & 500000\\\hline
L02 & C02 & Personal & 150000\\\hline
L03 & C03 & Business & 700000\\\hline
L04 & C04 & Personal & 120000\\
\hline

\end{tabular}
\chSixNextColumn
একজন গ্রাহকের একাধিক ব্যাংক হিসাব থাকতে পারে।

তাই ধারণাগতভাবে—

\[
Customer
\quad
1:M
\quad
Account
\]
\end{chSixSideBySide}
\noindent\begin{minipage}{\linewidth}
\chSixExample{৭.৪৯}{সব ব্যাংক হিসাবের তথ্য দেখাও।}

\begin{chSixSideBySide}[.32]
\begin{verbatim}
SELECT *
FROM Account;
\end{verbatim}
\chSixNextColumn
\textbf{ফল}\par\smallskip
{\raggedright\footnotesize\setlength{\tabcolsep}{2pt}
\begin{tabular}{|l|l|l|l|}\hline\chSixHeader
AccountNo & CustomerID & AccountType & Balance\\\hline
A1001 & C01 & Savings & 150000\\\hline
A1002 & C02 & Current & 85000\\\hline
A1003 & C01 & Current & 70000\\\hline
A1004 & C03 & Savings & 220000\\\hline
A1005 & C04 & Savings & 45000\\
\hline\end{tabular}\par}
\end{chSixSideBySide}
\end{minipage}\par\smallskip

\noindent\begin{minipage}{\linewidth}
\chSixExample{৭.৫০}{যেসব হিসাবের Balance 50000-এর বেশি, সেগুলো দেখাও।}

\noindent\begin{minipage}[c]{.46\linewidth}
\begin{verbatim}
SELECT AccountNo, CustomerID, Balance
FROM Account
WHERE Balance>50000;
\end{verbatim}
\end{minipage}\hfill\begin{minipage}[c]{.04\linewidth}\raggedright{\color{chSixTeal}$\Rightarrow$}\end{minipage}\hfill
\begin{minipage}[c]{.47\linewidth}\raggedright\footnotesize\setlength{\tabcolsep}{2pt}\textbf{ফল}\par\smallskip
\begin{tabular}{|l|l|l|}\hline\chSixHeader
AccountNo & CustomerID & Balance\\\hline
A1001 & C01 & 150000\\\hline
A1002 & C02 & 85000\\\hline
A1003 & C01 & 70000\\\hline
A1004 & C03 & 220000\\
\hline\end{tabular}
\end{minipage}
\end{minipage}\par\smallskip

\noindent\begin{minipage}{\linewidth}
\chSixExample{৭.৫১}{ঋণের পরিমাণের সর্বোচ্চ মান বের করো।}

\noindent\begin{minipage}[c]{.53\linewidth}
\begin{verbatim}
SELECT MAX(Amount) AS HighestLoan
FROM Loan;
\end{verbatim}
\end{minipage}\hfill\begin{minipage}[c]{.04\linewidth}\raggedright{\color{chSixTeal}$\Rightarrow$}\end{minipage}\hfill
\begin{minipage}[c]{.40\linewidth}\raggedright\footnotesize\setlength{\tabcolsep}{2pt}\textbf{ফল}\par\smallskip
\begin{tabular}{|l|}\hline\chSixHeader
HighestLoan\\\hline
700000\\
\hline\end{tabular}
\end{minipage}
\end{minipage}\par\smallskip

\noindent\begin{minipage}{\linewidth}
\chSixExample{৭.৫২}{Savings হিসাবগুলোর মধ্যে যেগুলোর Balance 100000-এর বেশি, সেগুলো দেখাও।}

\noindent\begin{minipage}[c]{.46\linewidth}
\begin{verbatim}
SELECT AccountNo, CustomerID, Balance
FROM Account
WHERE AccountType='Savings'
  AND Balance>100000;
\end{verbatim}
\end{minipage}\hfill\begin{minipage}[c]{.04\linewidth}\raggedright{\color{chSixTeal}$\Rightarrow$}\end{minipage}\hfill
\begin{minipage}[c]{.47\linewidth}\raggedright\footnotesize\setlength{\tabcolsep}{2pt}\textbf{ফল}\par\smallskip
\begin{tabular}{|l|l|l|}\hline\chSixHeader
AccountNo & CustomerID & Balance\\\hline
A1001 & C01 & 150000\\\hline
A1004 & C03 & 220000\\
\hline\end{tabular}
\end{minipage}
\end{minipage}\par\smallskip

ফলে A1001 ও A1004 আসবে।

\noindent\begin{minipage}{\linewidth}
\chSixExample{৭.৫৩}{যেসব হিসাবের Balance 80000 থেকে 180000-এর মধ্যে, সেগুলো দেখাও।}

\noindent\begin{minipage}[c]{.53\linewidth}
\begin{verbatim}
SELECT AccountNo, Balance
FROM Account
WHERE Balance BETWEEN 80000 AND 180000;
\end{verbatim}
\end{minipage}\hfill\begin{minipage}[c]{.04\linewidth}\raggedright{\color{chSixTeal}$\Rightarrow$}\end{minipage}\hfill
\begin{minipage}[c]{.40\linewidth}\raggedright\footnotesize\setlength{\tabcolsep}{2pt}\textbf{ফল}\par\smallskip
\begin{tabular}{|l|l|}\hline\chSixHeader
AccountNo & Balance\\\hline
A1001 & 150000\\\hline
A1002 & 85000\\
\hline\end{tabular}
\end{minipage}
\end{minipage}\par\smallskip

\noindent\begin{minipage}{\linewidth}
\chSixExample{৭.৫৪}{সব হিসাবের Balance যোগ করে মোট অর্থের পরিমাণ বের করো।}

\noindent\begin{minipage}[c]{.53\linewidth}
\begin{verbatim}
SELECT SUM(Balance) AS TotalBalance
FROM Account;
\end{verbatim}
\end{minipage}\hfill\begin{minipage}[c]{.04\linewidth}\raggedright{\color{chSixTeal}$\Rightarrow$}\end{minipage}\hfill
\begin{minipage}[c]{.40\linewidth}\raggedright\footnotesize\setlength{\tabcolsep}{2pt}\textbf{ফল}\par\smallskip
\begin{tabular}{|l|}\hline\chSixHeader
TotalBalance\\\hline
570000\\
\hline\end{tabular}
\end{minipage}
\end{minipage}\par\smallskip

\par TotalBalance = 150000+85000+70000+220000+45000 = 570000.\par

\noindent\begin{minipage}{\linewidth}
\chSixExample{৭.৫৫}{Balance-এর গড় বের করো।}

\noindent\begin{minipage}[c]{.53\linewidth}
\begin{verbatim}
SELECT AVG(Balance) AS AverageBalance
FROM Account;
\end{verbatim}
\end{minipage}\hfill\begin{minipage}[c]{.04\linewidth}\raggedright{\color{chSixTeal}$\Rightarrow$}\end{minipage}\hfill
\begin{minipage}[c]{.40\linewidth}\raggedright\footnotesize\setlength{\tabcolsep}{2pt}\textbf{ফল}\par\smallskip
\begin{tabular}{|l|}\hline\chSixHeader
AverageBalance\\\hline
114000\\
\hline\end{tabular}
\end{minipage}
\end{minipage}\par\smallskip

\par AverageBalance = 570000/5 = 114000.\par

\noindent\begin{minipage}{\linewidth}
\chSixExample{৭.৫৬}{Balance বড় থেকে ছোট ক্রমে ব্যাংক হিসাবগুলো সাজাও।}

\noindent\begin{minipage}[c]{.53\linewidth}
\begin{verbatim}
SELECT AccountNo, Balance
FROM Account
ORDER BY Balance DESC;
\end{verbatim}
\end{minipage}\hfill\begin{minipage}[c]{.04\linewidth}\raggedright{\color{chSixTeal}$\Rightarrow$}\end{minipage}\hfill
\begin{minipage}[c]{.40\linewidth}\raggedright\footnotesize\setlength{\tabcolsep}{2pt}\textbf{ফল}\par\smallskip
\chSixResultList{AccountNo, Balance}{\mbox{A1004, 220000}; \mbox{A1001, 150000}; \mbox{A1002, 85000}; \mbox{A1003, 70000}; \mbox{A1005, 45000}}
\end{minipage}
\end{minipage}\par\smallskip

\noindent\begin{minipage}{\linewidth}
\chSixExample{৭.৫৭}{A1002 হিসাবে 15000 টাকা জমা হলে তার Balance হালনাগাদ করো।}

\noindent\begin{minipage}[c]{.53\linewidth}
\begin{verbatim}
UPDATE Account
SET Balance=Balance+15000
WHERE AccountNo='A1002';
\end{verbatim}
\end{minipage}\hfill\begin{minipage}[c]{.04\linewidth}\raggedright{\color{chSixTeal}$\Rightarrow$}\end{minipage}\hfill
\begin{minipage}[c]{.40\linewidth}\raggedright\footnotesize\setlength{\tabcolsep}{2pt}\textbf{ফল}\par\smallskip
\begin{tabular}{|l|l|}\hline\chSixHeader
AccountNo & Balance\\\hline
A1002 & 100000\\
\hline\end{tabular}
\end{minipage}
\end{minipage}\par\smallskip

এই হিসাবে আগে 85000 টাকা ছিল। নতুন করে 15000 টাকা জমা হওয়ায় Balance হলো 85000+15000=100000 টাকা। তাই UPDATE-এ পুরোনো Balance-এর সঙ্গে জমার পরিমাণ যোগ করা হয়েছে।

\begin{keypointbox}{Calculated UPDATE}

\par\noindent\begin{minipage}{\linewidth}
\begin{verbatim}
SET Balance=Balance+15000
\end{verbatim}
\end{minipage}\par

এখানে পুরোনো Balance-এর সঙ্গে নতুন জমার পরিমাণ যোগ করা হচ্ছে।

\end{keypointbox}

